% --- Archon physics formalization source begin ---
% archon:physics
% archon:covers PhyXMiniProblems/problem_phyx_mini_0992.lean
% archon:source-report reports/phyx_mini/problem_phyx_mini_0992.source.json

\chapter{Physics problem phyx\_mini\_0992}
\label{ch:PhyXMiniProblems_problem_phyx_mini_0992}

\paragraph{Problem source.}
\#\# Physical scenario

Two point charges, \textbackslash{}( q\_1 = +25\textbackslash{},\textbackslash{}text\{nC\} \textbackslash{}) and \textbackslash{}( q\_2 = -75\textbackslash{},\textbackslash{}text\{nC\} \textbackslash{}), are separated by a distance \textbackslash{}( r = 3.0\textbackslash{},\textbackslash{}text\{cm\} \textbackslash{}).

\#\# Question

Find the magnitude of the electric force that \textbackslash{}( q\_1 \textbackslash{}) exerts on \textbackslash{}( q\_2 \textbackslash{}).

\#\# Answer choices

- A: A: 0.63N

- B: B: 0.0187N

- C: C: 0.019N

- D: D: 0.028N

\#\# Figure caption (auxiliary; use the image as primary evidence)

The image illustrates a basic representation of electrostatic charges. It features two spheres:

1. The sphere on the left is red with a plus sign, indicating a positive charge, labeled as \textbackslash{}( q\_1 \textbackslash{}).
2. The sphere on the right is blue with a minus sign, representing a negative charge, labeled as \textbackslash{}( q\_2 \textbackslash{}).

Both spheres are connected by a horizontal double-headed arrow labeled \textbackslash{}( r \textbackslash{}), indicating the distance between the two charges. The setup suggests an example related to Coulomb's law or electrostatic interactions.

\#\# Dataset metadata

Electrostatics; Physical Model Grounding Reasoning

\paragraph{Recorded answer/context.}
B: B: 0.0187N

\paragraph{Figure/image path.}
/root/proposal\_for\_physic/science-mango/phyx\_mini\_run/phyx\_data/test\_image/992.png

\paragraph{Formalization target.}
create a compiling Lean file with sorry bodies at `PhyXMiniProblems/problem\_phyx\_mini\_0992.lean`. The Lean declarations must preserve the physical quantities, dimensions or dimensional roles, figure labels, governing-law hypotheses, and final relation expressed by this problem.
Use Mathlib/Physlib names found through LeanExplore where available. If a physics API is missing, introduce faithful local abstractions rather than scalar placeholder aliases.

\begin{theorem}[Physics formalization target]
\label{thm:physics:phyx_mini_0992:target}
\lean{PhyXMiniProblems.ProblemPhyXMini0992.problem_phyx_mini_0992}
\uses{def:physics:phyx-mini-0992:phyxminiproblems-problemphyxmini0992-forcemagnitudeinnewtons, def:physics:phyx-mini-0992:phyxminiproblems-problemphyxmini0992-twopointchargeforcesetup, def:physics:phyx-mini-0992:phyxminiproblems-problemphyxmini0992-matchessuppliedtwopointchargefigure, def:physics:phyx-mini-0992:phyxminiproblems-problemphyxmini0992-matchesstatedpointchargedata, def:physics:phyx-mini-0992:phyxminiproblems-problemphyxmini0992-usesroundedcoulombconstant, def:physics:phyx-mini-0992:phyxminiproblems-problemphyxmini0992-hasphysicalpointchargeparameters, def:physics:phyx-mini-0992:phyxminiproblems-problemphyxmini0992-satisfiespointchargecoulombmagnitudelaw, def:physics:phyx-mini-0992:phyxminiproblems-problemphyxmini0992-recordeddatasetanswer, def:physics:phyx-mini-0992:phyxminiproblems-problemphyxmini0992-roundstofourdecimalplaces, def:physics:phyx-mini-0992:phyxminiproblems-problemphyxmini0992-isstrictlyclosestanswerchoice}
The assigned autoformalize agent should translate this physics problem into Lean declarations in the covered file, with theorem and lemma proof bodies written as `by sorry`.
\end{theorem}
\begin{proof}
This is an autoformalization task, not a proof task. Produce statements that can later be proved by the physics prover without weakening the physical model.
\end{proof}

% --- Archon declaration topology begin ---
\section{Declaration topology}
\begin{definition}[force Dimension]
  \label{def:physics:phyx-mini-0992:phyxminiproblems-problemphyxmini0992-forcedimension}
  \lean{PhyXMiniProblems.ProblemPhyXMini0992.forceDimension}
  The physical dimension M L T⁻² of force
\end{definition}
\begin{proof}
  This is the declared model definition of force Dimension.
\end{proof}
\begin{definition}[Signed Charge Quantity]
  \label{def:physics:phyx-mini-0992:phyxminiproblems-problemphyxmini0992-signedchargequantity}
  \lean{PhyXMiniProblems.ProblemPhyXMini0992.SignedChargeQuantity}
  A signed, unit-independent physical electric charge
\end{definition}
\begin{proof}
  This is the declared model definition of Signed Charge Quantity.
\end{proof}
\begin{definition}[Length Quantity]
  \label{def:physics:phyx-mini-0992:phyxminiproblems-problemphyxmini0992-lengthquantity}
  \lean{PhyXMiniProblems.ProblemPhyXMini0992.LengthQuantity}
  A nonnegative, unit-independent physical separation
\end{definition}
\begin{proof}
  This is the declared model definition of Length Quantity.
\end{proof}
\begin{definition}[Force Magnitude Quantity]
  \label{def:physics:phyx-mini-0992:phyxminiproblems-problemphyxmini0992-forcemagnitudequantity}
  \lean{PhyXMiniProblems.ProblemPhyXMini0992.ForceMagnitudeQuantity}
  \uses{def:physics:phyx-mini-0992:phyxminiproblems-problemphyxmini0992-forcedimension}
  A nonnegative, unit-independent physical force magnitude
\end{definition}
\begin{proof}
  This is the declared model definition of Force Magnitude Quantity.
\end{proof}
\begin{definition}[charge In Coulombs]
  \label{def:physics:phyx-mini-0992:phyxminiproblems-problemphyxmini0992-chargeincoulombs}
  \lean{PhyXMiniProblems.ProblemPhyXMini0992.chargeInCoulombs}
  \uses{def:physics:phyx-mini-0992:phyxminiproblems-problemphyxmini0992-signedchargequantity}
  Coherent-SI readout of a signed charge, in coulombs
\end{definition}
\begin{proof}
  This is the declared model definition of charge In Coulombs.
\end{proof}
\begin{definition}[charge In Nanocoulombs]
  \label{def:physics:phyx-mini-0992:phyxminiproblems-problemphyxmini0992-chargeinnanocoulombs}
  \lean{PhyXMiniProblems.ProblemPhyXMini0992.chargeInNanocoulombs}
  \uses{def:physics:phyx-mini-0992:phyxminiproblems-problemphyxmini0992-signedchargequantity, def:physics:phyx-mini-0992:phyxminiproblems-problemphyxmini0992-chargeincoulombs}
  Nanocoulomb readout used by the numerical problem statement
\end{definition}
\begin{proof}
  This is the declared model definition of charge In Nanocoulombs.
\end{proof}
\begin{definition}[length In Meters]
  \label{def:physics:phyx-mini-0992:phyxminiproblems-problemphyxmini0992-lengthinmeters}
  \lean{PhyXMiniProblems.ProblemPhyXMini0992.lengthInMeters}
  \uses{def:physics:phyx-mini-0992:phyxminiproblems-problemphyxmini0992-lengthquantity}
  Coherent-SI readout of a separation, in metres
\end{definition}
\begin{proof}
  This is the declared model definition of length In Meters.
\end{proof}
\begin{definition}[length In Centimeters]
  \label{def:physics:phyx-mini-0992:phyxminiproblems-problemphyxmini0992-lengthincentimeters}
  \lean{PhyXMiniProblems.ProblemPhyXMini0992.lengthInCentimeters}
  \uses{def:physics:phyx-mini-0992:phyxminiproblems-problemphyxmini0992-lengthquantity, def:physics:phyx-mini-0992:phyxminiproblems-problemphyxmini0992-lengthinmeters}
  Centimetre readout used by the numerical problem statement
\end{definition}
\begin{proof}
  This is the declared model definition of length In Centimeters.
\end{proof}
\begin{definition}[force Magnitude In Newtons]
  \label{def:physics:phyx-mini-0992:phyxminiproblems-problemphyxmini0992-forcemagnitudeinnewtons}
  \lean{PhyXMiniProblems.ProblemPhyXMini0992.forceMagnitudeInNewtons}
  \uses{def:physics:phyx-mini-0992:phyxminiproblems-problemphyxmini0992-forcemagnitudequantity}
  Coherent-SI readout of a force magnitude, in newtons
\end{definition}
\begin{proof}
  This is the declared model definition of force Magnitude In Newtons.
\end{proof}
\begin{definition}[Charge Label]
  \label{def:physics:phyx-mini-0992:phyxminiproblems-problemphyxmini0992-chargelabel}
  \lean{PhyXMiniProblems.ProblemPhyXMini0992.ChargeLabel}
  The two charges named in image 992.png and in the problem text
\end{definition}
\begin{proof}
  This is the declared model definition of Charge Label.
\end{proof}
\begin{definition}[Horizontal Side]
  \label{def:physics:phyx-mini-0992:phyxminiproblems-problemphyxmini0992-horizontalside}
  \lean{PhyXMiniProblems.ProblemPhyXMini0992.HorizontalSide}
  Horizontal location of a charge marker in the supplied image
\end{definition}
\begin{proof}
  This is the declared model definition of Horizontal Side.
\end{proof}
\begin{definition}[Charge Sphere Color]
  \label{def:physics:phyx-mini-0992:phyxminiproblems-problemphyxmini0992-chargespherecolor}
  \lean{PhyXMiniProblems.ProblemPhyXMini0992.ChargeSphereColor}
  Fill color of a charge sphere in the primary image
\end{definition}
\begin{proof}
  This is the declared model definition of Charge Sphere Color.
\end{proof}
\begin{definition}[Charge Sign Glyph]
  \label{def:physics:phyx-mini-0992:phyxminiproblems-problemphyxmini0992-chargesignglyph}
  \lean{PhyXMiniProblems.ProblemPhyXMini0992.ChargeSignGlyph}
  Sign glyph drawn inside a charge sphere
\end{definition}
\begin{proof}
  This is the declared model definition of Charge Sign Glyph.
\end{proof}
\begin{definition}[Charged Body Model]
  \label{def:physics:phyx-mini-0992:phyxminiproblems-problemphyxmini0992-chargedbodymodel}
  \lean{PhyXMiniProblems.ProblemPhyXMini0992.ChargedBodyModel}
  Idealization of each charged body used in the governing law
\end{definition}
\begin{proof}
  This is the declared model definition of Charged Body Model.
\end{proof}
\begin{definition}[expected Charge Name]
  \label{def:physics:phyx-mini-0992:phyxminiproblems-problemphyxmini0992-expectedchargename}
  \lean{PhyXMiniProblems.ProblemPhyXMini0992.expectedChargeName}
  \uses{def:physics:phyx-mini-0992:phyxminiproblems-problemphyxmini0992-chargelabel}
  Expected typographic label beside each charge sphere
\end{definition}
\begin{proof}
  This is the declared model definition of expected Charge Name.
\end{proof}
\begin{definition}[expected Horizontal Side]
  \label{def:physics:phyx-mini-0992:phyxminiproblems-problemphyxmini0992-expectedhorizontalside}
  \lean{PhyXMiniProblems.ProblemPhyXMini0992.expectedHorizontalSide}
  \uses{def:physics:phyx-mini-0992:phyxminiproblems-problemphyxmini0992-chargelabel, def:physics:phyx-mini-0992:phyxminiproblems-problemphyxmini0992-horizontalside}
  Expected left/right placement of the two charge spheres
\end{definition}
\begin{proof}
  This is the declared model definition of expected Horizontal Side.
\end{proof}
\begin{definition}[expected Sphere Color]
  \label{def:physics:phyx-mini-0992:phyxminiproblems-problemphyxmini0992-expectedspherecolor}
  \lean{PhyXMiniProblems.ProblemPhyXMini0992.expectedSphereColor}
  \uses{def:physics:phyx-mini-0992:phyxminiproblems-problemphyxmini0992-chargelabel, def:physics:phyx-mini-0992:phyxminiproblems-problemphyxmini0992-chargespherecolor}
  Expected sphere color at each labelled charge
\end{definition}
\begin{proof}
  This is the declared model definition of expected Sphere Color.
\end{proof}
\begin{definition}[expected Charge Sign]
  \label{def:physics:phyx-mini-0992:phyxminiproblems-problemphyxmini0992-expectedchargesign}
  \lean{PhyXMiniProblems.ProblemPhyXMini0992.expectedChargeSign}
  \uses{def:physics:phyx-mini-0992:phyxminiproblems-problemphyxmini0992-chargelabel, def:physics:phyx-mini-0992:phyxminiproblems-problemphyxmini0992-chargesignglyph}
  Expected sign glyph at each labelled charge
\end{definition}
\begin{proof}
  This is the declared model definition of expected Charge Sign.
\end{proof}
\begin{definition}[Two Point Charge Figure]
  \label{def:physics:phyx-mini-0992:phyxminiproblems-problemphyxmini0992-twopointchargefigure}
  \lean{PhyXMiniProblems.ProblemPhyXMini0992.TwoPointChargeFigure}
  \uses{def:physics:phyx-mini-0992:phyxminiproblems-problemphyxmini0992-chargelabel, def:physics:phyx-mini-0992:phyxminiproblems-problemphyxmini0992-horizontalside, def:physics:phyx-mini-0992:phyxminiproblems-problemphyxmini0992-chargespherecolor, def:physics:phyx-mini-0992:phyxminiproblems-problemphyxmini0992-chargesignglyph}
  Literal presentation data transcribed from image 992.png. It contains no force value and no answer-choice label
\end{definition}
\begin{proof}
  This is the declared model definition of Two Point Charge Figure.
\end{proof}
\begin{definition}[Two Point Charge Force Setup]
  \label{def:physics:phyx-mini-0992:phyxminiproblems-problemphyxmini0992-twopointchargeforcesetup}
  \lean{PhyXMiniProblems.ProblemPhyXMini0992.TwoPointChargeForceSetup}
  \uses{def:physics:phyx-mini-0992:phyxminiproblems-problemphyxmini0992-signedchargequantity, def:physics:phyx-mini-0992:phyxminiproblems-problemphyxmini0992-lengthquantity, def:physics:phyx-mini-0992:phyxminiproblems-problemphyxmini0992-forcemagnitudequantity, def:physics:phyx-mini-0992:phyxminiproblems-problemphyxmini0992-chargelabel, def:physics:phyx-mini-0992:phyxminiproblems-problemphyxmini0992-chargedbodymodel, def:physics:phyx-mini-0992:phyxminiproblems-problemphyxmini0992-twopointchargefigure}
  Independent physical quantities in the two-charge setup. In particular, the force exerted by q₁ on q₂ is stored independently rather than defined from Coulomb's formula or from the recorded answer
\end{definition}
\begin{proof}
  This is the declared model definition of Two Point Charge Force Setup.
\end{proof}
\begin{definition}[Matches Supplied Two Point Charge Figure]
  \label{def:physics:phyx-mini-0992:phyxminiproblems-problemphyxmini0992-matchessuppliedtwopointchargefigure}
  \lean{PhyXMiniProblems.ProblemPhyXMini0992.MatchesSuppliedTwoPointChargeFigure}
  \uses{def:physics:phyx-mini-0992:phyxminiproblems-problemphyxmini0992-expectedchargename, def:physics:phyx-mini-0992:phyxminiproblems-problemphyxmini0992-expectedhorizontalside, def:physics:phyx-mini-0992:phyxminiproblems-problemphyxmini0992-expectedspherecolor, def:physics:phyx-mini-0992:phyxminiproblems-problemphyxmini0992-expectedchargesign, def:physics:phyx-mini-0992:phyxminiproblems-problemphyxmini0992-twopointchargeforcesetup}
  Exact qualitative evidence read from the supplied primary image
\end{definition}
\begin{proof}
  This is the declared model definition of Matches Supplied Two Point Charge Figure.
\end{proof}
\begin{definition}[Matches Stated Point Charge Data]
  \label{def:physics:phyx-mini-0992:phyxminiproblems-problemphyxmini0992-matchesstatedpointchargedata}
  \lean{PhyXMiniProblems.ProblemPhyXMini0992.MatchesStatedPointChargeData}
  \uses{def:physics:phyx-mini-0992:phyxminiproblems-problemphyxmini0992-chargeinnanocoulombs, def:physics:phyx-mini-0992:phyxminiproblems-problemphyxmini0992-lengthincentimeters, def:physics:phyx-mini-0992:phyxminiproblems-problemphyxmini0992-twopointchargeforcesetup}
  The numerical values and point-charge idealization stated in the prose. These are calibrated readouts of dimensionful quantities, not scalar replacements for charge or length
\end{definition}
\begin{proof}
  This is the declared model definition of Matches Stated Point Charge Data.
\end{proof}
\begin{definition}[Uses Rounded Coulomb Constant]
  \label{def:physics:phyx-mini-0992:phyxminiproblems-problemphyxmini0992-usesroundedcoulombconstant}
  \lean{PhyXMiniProblems.ProblemPhyXMini0992.UsesRoundedCoulombConstant}
  \uses{def:physics:phyx-mini-0992:phyxminiproblems-problemphyxmini0992-twopointchargeforcesetup}
  The rounded textbook calibration 8.99 × 10⁹ N m²/C². Physlib represents Electromagnetism.EMSystem.coulombConstant as the coherent-SI scalar 1 / (4 π ε₀); this premise supplies the precision used by the answer list
\end{definition}
\begin{proof}
  This is the declared model definition of Uses Rounded Coulomb Constant.
\end{proof}
\begin{definition}[Has Physical Point Charge Parameters]
  \label{def:physics:phyx-mini-0992:phyxminiproblems-problemphyxmini0992-hasphysicalpointchargeparameters}
  \lean{PhyXMiniProblems.ProblemPhyXMini0992.HasPhysicalPointChargeParameters}
  \uses{def:physics:phyx-mini-0992:phyxminiproblems-problemphyxmini0992-lengthinmeters, def:physics:phyx-mini-0992:phyxminiproblems-problemphyxmini0992-twopointchargeforcesetup}
  Positivity and separation conditions selecting the physical branch
\end{definition}
\begin{proof}
  This is the declared model definition of Has Physical Point Charge Parameters.
\end{proof}
\begin{definition}[Satisfies Point Charge Coulomb Magnitude Law]
  \label{def:physics:phyx-mini-0992:phyxminiproblems-problemphyxmini0992-satisfiespointchargecoulombmagnitudelaw}
  \lean{PhyXMiniProblems.ProblemPhyXMini0992.SatisfiesPointChargeCoulombMagnitudeLaw}
  \uses{def:physics:phyx-mini-0992:phyxminiproblems-problemphyxmini0992-chargeincoulombs, def:physics:phyx-mini-0992:phyxminiproblems-problemphyxmini0992-lengthinmeters, def:physics:phyx-mini-0992:phyxminiproblems-problemphyxmini0992-forcemagnitudeinnewtons, def:physics:phyx-mini-0992:phyxminiproblems-problemphyxmini0992-twopointchargeforcesetup}
  Coulomb's inverse-square law for the magnitude of the force exerted by q₁ on q₂. Absolute value retains the signed charge data while producing a nonnegative magnitude. This general law contains no numerical target value
\end{definition}
\begin{proof}
  This is the declared model definition of Satisfies Point Charge Coulomb Magnitude Law.
\end{proof}
\begin{definition}[Answer Choice]
  \label{def:physics:phyx-mini-0992:phyxminiproblems-problemphyxmini0992-answerchoice}
  \lean{PhyXMiniProblems.ProblemPhyXMini0992.AnswerChoice}
  The four force magnitudes printed in the source's answer list
\end{definition}
\begin{proof}
  This is the declared model definition of Answer Choice.
\end{proof}
\begin{definition}[force In Newtons]
  \label{def:physics:phyx-mini-0992:phyxminiproblems-problemphyxmini0992-answerchoice-forceinnewtons}
  \lean{PhyXMiniProblems.ProblemPhyXMini0992.AnswerChoice.forceInNewtons}
  \uses{def:physics:phyx-mini-0992:phyxminiproblems-problemphyxmini0992-answerchoice}
  Literal newton value printed beside an answer choice
\end{definition}
\begin{proof}
  This is the declared model definition of force In Newtons.
\end{proof}
\begin{definition}[recorded Dataset Answer]
  \label{def:physics:phyx-mini-0992:phyxminiproblems-problemphyxmini0992-recordeddatasetanswer}
  \lean{PhyXMiniProblems.ProblemPhyXMini0992.recordedDatasetAnswer}
  \uses{def:physics:phyx-mini-0992:phyxminiproblems-problemphyxmini0992-answerchoice}
  The answer label recorded in the supplied dataset metadata
\end{definition}
\begin{proof}
  This is the declared model definition of recorded Dataset Answer.
\end{proof}
\begin{definition}[Rounds To Four Decimal Places]
  \label{def:physics:phyx-mini-0992:phyxminiproblems-problemphyxmini0992-roundstofourdecimalplaces}
  \lean{PhyXMiniProblems.ProblemPhyXMini0992.RoundsToFourDecimalPlaces}
  The strict half-unit criterion for displaying an actual newton value to four places after the decimal point. The strict inequality excludes tie cases
\end{definition}
\begin{proof}
  This is the declared model definition of Rounds To Four Decimal Places.
\end{proof}
\begin{definition}[Is Strictly Closest Answer Choice]
  \label{def:physics:phyx-mini-0992:phyxminiproblems-problemphyxmini0992-isstrictlyclosestanswerchoice}
  \lean{PhyXMiniProblems.ProblemPhyXMini0992.IsStrictlyClosestAnswerChoice}
  \uses{def:physics:phyx-mini-0992:phyxminiproblems-problemphyxmini0992-answerchoice}
  A displayed answer is uniquely closest to the actual force magnitude
\end{definition}
\begin{proof}
  This is the declared model definition of Is Strictly Closest Answer Choice.
\end{proof}
\begin{lemma}[force On Q2 From Q1 eq coulomb Expression]
  \label{lem:physics:phyx-mini-0992:phyxminiproblems-problemphyxmini0992-forceonq2fromq1-eq-coulombexpression}
  \lean{PhyXMiniProblems.ProblemPhyXMini0992.forceOnQ2FromQ1_eq_coulombExpression}
  \uses{def:physics:phyx-mini-0992:phyxminiproblems-problemphyxmini0992-chargeincoulombs, def:physics:phyx-mini-0992:phyxminiproblems-problemphyxmini0992-lengthinmeters, def:physics:phyx-mini-0992:phyxminiproblems-problemphyxmini0992-forcemagnitudeinnewtons, def:physics:phyx-mini-0992:phyxminiproblems-problemphyxmini0992-twopointchargeforcesetup, def:physics:phyx-mini-0992:phyxminiproblems-problemphyxmini0992-hasphysicalpointchargeparameters, def:physics:phyx-mini-0992:phyxminiproblems-problemphyxmini0992-satisfiespointchargecoulombmagnitudelaw}
  Specializing the governing magnitude law gives the exact coherent-SI Coulomb expression before any decimal rounding
\end{lemma}
\begin{proof}
  The conclusion follows from the governing relations and figure readouts named in the statement of force On Q2 From Q1 eq coulomb Expression.
\end{proof}
\begin{lemma}[stated Data In Coherent SI]
  \label{lem:physics:phyx-mini-0992:phyxminiproblems-problemphyxmini0992-stateddataincoherentsi}
  \lean{PhyXMiniProblems.ProblemPhyXMini0992.statedDataInCoherentSI}
  \uses{def:physics:phyx-mini-0992:phyxminiproblems-problemphyxmini0992-chargeincoulombs, def:physics:phyx-mini-0992:phyxminiproblems-problemphyxmini0992-lengthinmeters, def:physics:phyx-mini-0992:phyxminiproblems-problemphyxmini0992-twopointchargeforcesetup, def:physics:phyx-mini-0992:phyxminiproblems-problemphyxmini0992-matchesstatedpointchargedata}
  The stated nanocoulomb and centimetre values have these coherent-SI readouts. This helper records unit conversion only; it does not mention the force
\end{lemma}
\begin{proof}
  The conclusion follows from the governing relations and figure readouts named in the statement of stated Data In Coherent SI.
\end{proof}
% --- Archon declaration topology end ---
% --- Archon physics formalization source end ---
